\documentclass[12pt,italian]{article}
\usepackage[italian]{babel}
\usepackage[a4paper, margin=1.97cm]{geometry}
\usepackage{graphicx}
\graphicspath{{./images/}}
\usepackage{amsmath}
\usepackage[american]{circuitikz}
\usepackage{caption}
\usepackage{subcaption}
\numberwithin{table}{section}
\usepackage{amsfonts}
\usepackage[bookmarks,pdfusetitle,hidelinks]{hyperref}
\makeatletter
\let\oldsection\section
\renewcommand{\section}{\@ifstar{\@mystarsection}{\@mynormsection}}
\newcommand{\@mystarsection}[1]{%
    \oldsection*{#1}%
    \phantomsection%
    \addcontentsline{toc}{section}{#1}%
}
\newcommand{\@mynormsection}[2][]{%
    \if\relax\detokenize{#1}\relax
        \oldsection{#2}%
    \else
        \oldsection[#1]{#2}%
    \fi
}
\makeatother
\usepackage{cleveref}
\usepackage{siunitx}
\usepackage{booktabs}

\newcommand{\err}[1]{\textcolor{red}{#1}}


\title{Misure su un circuito tosatore a due livelli.}
\author{Enrico Barbuio \\ 0001117553 \and Giacomo Cicala \\ 0001122965} 
\date{Turno 1 del 12 Marzo 2026}

\begin{document}

\maketitle

\begin{abstract}

\end{abstract}
\section*{Introduzione}

\section*{Apparato sperimentale e svolgimento}
Il materiale utilizzato nella prova consiste in:\\
Strumenti:
\begin{itemize}
	\item Generatore di tensione costante (modello \texttt{TTi EB2025T}) con uscite
	      impostate a \qty{+5}{\volt} e \qty{-5}{\volt}.
	\item Generatore di funzioni (\err{???}).
	\item Oscilloscopio (\err{???}).
	\item Multimetro \texttt{Fluke 175 True RMS Multimeter}.
\end{itemize}
Componenti:
\begin{itemize}
	\item Potenziometro \qty{50}{\kilo\ohm} ($R$)
	\item 2x Potenziometri \qty{1}{\kilo\ohm} ($R_1$ e $R_2$)
	\item 2x Diodi al silicio / 2x Diodi al germanio ($D_1$ e $D_2$).
\end{itemize}

\begin{figure}[h]
	\centering
	\begin{circuitikz}
		\coordinate (start) at (0,0);
		\node[ground] at (start) {};

		\draw (start) to[sinusoidal voltage source, l={$V_{in}$}] ++(0, 2)
         to[variable resistor, l={$R$}] ++(0, 2) coordinate (top_left);

		\draw (top_left) -- ++(11, 0) coordinate (top_right);
		\draw (start) -- ++(11, 0) coordinate (bot_right);
		\draw (bot_right) to[oscope, l_={$V_{out}$}] (top_right);

        % tosatore 1
		\draw (start) ++(3, 1.5) coordinate (R1_left) to[battery2, v_=\qty{5}{\volt}] ++(0, -1.5);

		\draw (R1_left) to[pR, name=pot1, l_={$R_1$}] ++(2, 0) coordinate (R1_right);
		\draw (R1_right) -| ++(0, -1.5);

		\draw (pot1.wiper) to[short, -o] ++(0.75,0) node[right]{$V_{R1}$};
		\draw (pot1.wiper |- top_left) to[full diode, l=$D_1$] (pot1.wiper);

        % tosatore 2
		\draw (start) ++(6, 0) -| ++(0, 1.5) coordinate (R2_left);
		\draw (R2_left) to[pR, name=pot2, l_={$R_2$}] ++(2, 0) coordinate (R2_right);

		\draw (start) ++(8, 0) to[battery2, v_=\qty{-5}{\volt}] (R2_right);

		\draw (pot2.wiper) to[short, -o] ++(0.75,0) node[right]{$V_{R2}$};
		\draw (pot2.wiper) to[full diode, l_=$D_2$] (pot2.wiper |- top_left);
	\end{circuitikz}
	\caption{Schema del circuito realizzato.}
	\label{fig:schema_circuito}
\end{figure}

Descriviamo la realizzazione del circuito in \cref{fig:schema_circuito}.
Il potenziometro $R$ è stato usato come una resistenza utilizzando due dei tre
piedini, è stato impostato a $(40.0 \pm 0.4)$ \unit{\kilo\ohm}.
Successivamente, variando i potenziometri $R_1$ e $R_2$ sono state impostate le
tensioni $V_{R1}$ e $V_{R2}$ a: $V_{R1} = (2.004 \pm 0.005)$ \unit{\volt} e 
 $V_{R2} = (-2.004 \pm 0.005)$ \unit{\volt}.

Dopodiché sono stati collegati l'oscilloscopio e il generatore di funzioni.
Il generatore di funzioni è stato impostato con una forma sinusoidale,
frequenza di \qty{1}{\kilo\hertz} e semiampiezza \qty{6}{\volt}.
Per conferma questi parametri sono stati misurati con l'oscilloscopio ottenendo:
\[
    V_{in, 0} = (6.0 \pm 0.1) \ \unit{\volt},
    \qquad \tau = (1.00 \pm 0.04) \ \unit{\milli\second} \quad \Rightarrow \quad
    \nu = (1.00 \pm 0.04 ) \ \unit{\kilo\hertz}.
\]
A questo punto è stata osservata sull'oscilloscopio la forma dell'onda tosata e
sono state svolte le misure dei tempi di salita e discesa.\err{descrivere?}

L'esperienza è stata ripetuta, prima con due diodi al silicio e poi con due
diodi al germanio.

Rimanendo poi con i diodi al silicio sono state fatte altre due prove.
Prima è stata abbassata la tensione $V_{R1}$, responsabile del clipping
superiore, a $(1.002 \pm 0.004)$ \unit{\volt}.
Infine, riportando $V_{R1}$ al valore iniziale, è stata misurata la resistenza
critica $R_c$,
ovvero il valore di $R$ per cui si inizia ad osservare una diminuzione
dell'effetto di taglio sull'onda sinusoidale.
Come valore di soglia è stata considerata una variazione del picco di una
tacchetta dell'oscilloscopio nella scala \err{??? \unit {\volt}}.

\section*{Risultati e discussione}
\section*{Conclusioni}

\appendix
\section*{Errori per le misure con l'oscilloscopio}
\err{non penso serva una sezione ma per ora li metto qua.}

Assumiamo innanzitutto di apprezzare la mezza tacca dell'oscilloscopio,
l'errore su ciascuna lettura sarà quindi $\Delta_i = \frac{{\text{F.S.}}_i}{2}$.
Le misure svolte con l'oscilloscopio sono sempre delle delta tra due valori per
cui l'incertezza totale è $\Delta = \Delta_1 + \Delta_2$.
Quando entrambe le misure vengono effettuate sulla stessa scala risulta
$\Delta = 2 \cdot \frac{\text{F.S.}}{2} = \text{F.S.}$.
Invece, nel caso di misure di tensione in riferimento alla massa, è stato
azzerato l'oscilloscopio con fondoscala di \qty{0.04}{\volt}, questo errore è
sempre ordini di grandezza inferiore all'errore sull'altra misura per cui:
$\Delta = \Delta_{GND} + \Delta_{val} \approx \Delta_{val}$.
\end{document}